Read in Gauge Theory Gravity, the Christoffel symbols of clause (iii) below are those of the covariant derivative $D$ in the frame of the stated gauge, and $\nabla_{\partial_v}\partial_v = \kappa\,\partial_v$ is $(K \cdot D)K = \kappa K$ (Fact \claimref{synthesis.imported-gtg-dictionary}, clauses (ii) and (vi)). Conventions. $a(f) = \int \overline{f(u)}\,\psi(u)\,du$ (antilinear in $f$), $a^{*}(f) = a(f)^{*}$; a gauge-invariant quasi-free state has symbol $A$, $\omega(a^{*}(f)a(g)) = \ip{g}{Af}$, so that its kernel is $A(u, v) = \omega(\psi^{*}(v)\psi(u))$; a flow $\tau$ is written $\tau_{t}(a(f)) = a(e^{I_{4}th}f)$ (the unit of the one-particle space is $I_{4}$; numbers are read in $\mathrm{span}\{1, I_{4}\}$), the convention of Theorem \claimref{hawking.i-state-continuum}, in which the I-state at $\beta$ has symbol $(1 + e^{\beta h})^{-1}$. Since $\nu_{a}(a(f)) = \int\overline{f(u)}\,\psi(u + a)\,du = a(f(\cdot - a))$, the one-particle generator of $\nu$ is $P = I_{4}\,d/du$ ($e^{I_{4}aP}f = f(\cdot - a)$). Since $\beta_{t}(a(g)) = \int\overline{g(u)}\,e^{\lambda/2}\psi(e^{\lambda}u)\,du = a\bigl(e^{-\lambda/2}g(e^{-\lambda}\,\cdot)\bigr)$ with $\lambda = \kappa t$, on $\varphi_{w}(u) = u^{-1/2 + I_{4}w}$ ($u > 0$) one has $e^{-\lambda/2}\varphi_{w}(e^{-\lambda}u) = e^{-I_{4}w\lambda}\varphi_{w}(u)$, so the one-particle boost generator $h$ acts on $\varphi_{w}$ as $-\kappa w$; the half-density factor $e^{\lambda/2}$ is what makes $\beta_{t}$ preserve $\{\psi(u), \psi^{*}(v)\} = \delta(u - v)$, since $\delta(e^{\lambda}(u - v)) = e^{-\lambda}\delta(u - v)$. The transverse factor $K$ is carried along as $\otimes\id_{K}$ throughout and is suppressed below.

\emph{(i)} Hypothesis (N$_{+}$) is Premise \claimref{synthesis.n-plus-null-ground-state}: $\omega$ is a gauge-invariant quasi-free state, invariant under $\nu$ with positive generator in its GNS representation, that is a gauge-invariant quasi-free ground state for $\nu$. By Theorem \claimref{synthesis.null-vacuum-on-a-horizon-generator}(ii), where the argument (formerly given here) now stands, such a state has symbol $A = \chi_{(-\infty,0)}(P)$, $P = I_{4}\,d/du$: invariance makes $A$ a function of $P$, and the ground-state characterisation (Bratteli and Robinson, Vol.~2, Definition 5.3.18 and Proposition 5.3.19, cited from memory) applied to $a(g)$ and $a^{*}(f)$ gives $PA \le 0$ and $(1 - A)P \ge 0$, which force $A(k) = 1$ for $k < 0$ and $A(k) = 0$ for $k > 0$. So $\omega = \omega_{0}$, the null vacuum, with kernel $\tfrac{I_{4}}{2\pi}(u - v + I_{4}0)^{-1}$, and no other gauge-invariant quasi-free state satisfies (N$_{+}$); conversely $\omega_{0}$ satisfies it (the same Theorem, (ii)).

\emph{(ii)} By (i) the state is $\omega_{0}$, and (ii) is Theorem \claimref{synthesis.null-vacuum-on-a-horizon-generator}(iii), where its argument (formerly given here) now stands: the modes beyond the edge form the $\beta$-invariant subalgebra $\CAR(L^{2}(\R_{+}, du))$; in the Mellin basis $\varphi_{w}(u) = u^{-1/2 + I_{4}w}$, on which the one-particle boost generator acts as $-\kappa w$, the compressed sea projection has symbol $m(w) = (1 + e^{-2\pi w})^{-1}$, which is $(1 + e^{\beta h})^{-1}$ exactly at $\beta = 2\pi/\kappa$, so by Theorem \claimref{hawking.i-state-continuum}(1) and Fact \claimref{hawking.imported-quasi-free-kms-reduces-two} the restriction is the continuum I-state, jointly KMS and even, and by Fact \claimref{hawking.imported-ideal-fermi-gas-uniqueness} the unique $\beta$-KMS state there; the reflection $u \mapsto -u$ gives the past half for the boost future-directed there, $\beta_{-t}$; positivity fixes the sign of the period, and the one-sided support is needed (on the whole generator the boost-continued two-point function has a pole at imaginary rapidity $\pi$). The local-horizon premise and the Einstein derivation built on it (Theorem \claimref{synthesis.einstein-equations-from-c-loc-and-clausius}) use the past half, generators ending on the edge, with future-directed boost $\chi = -\xi$. The checks of this result still verify the strip, the past half and the Mellin symbol at sample points.

\emph{(iii)} The horizons. For (a), in ingoing Eddington-Finkelstein form $f\,dv^{2} - 2\,dv\,dr - r^{2}d\Omega^{2}$ (store signature) with $f = 1 - 2M/r$, the generator of the future horizon $r = 2M$ is the physical Killing vector field $K = \hb^{-1}(e_{v})$ over $\partial_{v}$, and on the horizon locus $(K \cdot D)K = \kappa K$ with $\kappa = 1/(4M)$, the value of Proposition \claimref{singularity.euclidean-section-period}: the component of $(K \cdot D)K$ off $K$ vanishes there (the generator stays on the horizon) and its component along $K$ is $\tfrac12 f'(2M)$ (in components, by Fact \claimref{synthesis.imported-gtg-dictionary}(ii): $\Gamma^{r}_{vv} = -\tfrac12 g^{rr}f' = \tfrac12 ff' = 0$ and $\Gamma^{v}_{vv} = -\tfrac12 g^{vr}f' = \tfrac12 f'(2M) = 1/(4M)$). For (b), the future cosmological horizon of the static patch is reached as $t \to +\infty$ at finite outgoing time $u_{\mathrm{out}} = t - r_{*}$, so it is computed in the outgoing form $f\,du_{\mathrm{out}}^{2} + 2\,du_{\mathrm{out}}\,dr - r^{2}d\Omega^{2}$ with $f = 1 - H^{2}r^{2}$: there the Killing field $K = \hb^{-1}(e_{u})$ has $(K \cdot D)K = HK$ on $r = 1/H$ (in components $g^{ur} = 1$, $\Gamma^{r}_{uu} = \tfrac12 ff' = 0$ and $\Gamma^{u}_{uu} = -\tfrac12 f'(1/H) = H$), the surface gravity $\kappa_{\Lambda} = H$ of Proposition \claimref{synthesis.de-sitter-static-boost-frozen-point}(iii). (The ingoing chart gives $(K \cdot D)K = -HK$ on $r = 1/H$, in components $\Gamma^{v}_{vv} = -H$: it covers the past cosmological horizon, where the future-directed Killing flow moves points towards the bifurcation sphere, the past-half case of (ii), with the same $|\kappa|$.) In both cases, with $s$ the Killing-time coordinate along the generator ($v$, resp. $u_{\mathrm{out}}$) and $\kappa > 0$, the affine parameter is $V = e^{\kappa s}/\kappa$ ($V'' = \kappa V'$), with $V \to 0$ at the bifurcation sphere as $s \to -\infty$; the generator continues to $V < 0$ in the Kruskal, respectively de Sitter, extension, where it bounds the other exterior, so (N$_{+}$) in cases (a) and (b) is posed on the whole generator $V \in \R$ of the extended solution, a Hartle-Hawking-type condition (Bunch-Davies-type in de Sitter). For a hole formed by collapse the generators of the event horizon begin at finite $V$, there is no $V < 0$ half, and the physical state is not a ground state for the affine translations on the horizon, so case (a) is the eternal hole only. The Killing field is $\partial_{s} = \kappa V\,\partial_{V}$: Killing time $t$ acts on each generator as $V \mapsto e^{\kappa t}V$, the boost $\beta_{t}$ with $u = V$. For (c), in the store signature the rotor $B = \cosh\tfrac{\lambda}{2} + \sinh\tfrac{\lambda}{2}\,e_{1}e_{0}$ gives $B\ell\tilde B = e^{\lambda}\ell$ for $\ell = e_{0} + e_{1}$, $Bn\tilde B = e^{-\lambda}n$ for $n = e_{0} - e_{1}$, and fixes $e_{2}, e_{3}$, so on the null plane $\{u\ell + y\}$ it is $u \mapsto e^{\lambda}u$ at fixed $y$: every generator dilated about its own edge point at unit rate in rapidity. For a gravitational field of nonvanishing field strength this holds at leading order, with the leading-order rapidity normalisation of Theorem \claimref{synthesis.local-boost-temperature-is-local-regularity}(i). The check verifies the non-affinities from the inverse metrics over $\Q(r)$ (ingoing chart at three masses; outgoing chart, $+H$, and ingoing chart, $-H$, at three values of $H$), the series identities $V'' = \kappa V'$ and $V' = \kappa V + 1$ for $V = (e^{\kappa v} - 1)/\kappa$ (the same Killing field about $V = -1/\kappa$) to order 12, and the rotor action at $e^{\lambda/2} = \tfrac32, \tfrac15, 7$. Then (ii) with these $\kappa$ gives: on (a) the I-state at $\beta = 8\pi M$, Fermi-Dirac in Killing frequency $\varpi = -\kappa w$, the continuum restatement of (R) for the eternal hole; on (b) the same at $2\pi/H$, that of (R$_{\Lambda}$); on (c), on the past half, KMS for the future-directed local boost at rapidity period $2\pi$ on the patch's whole mode algebra, that of (R$_{\mathrm{loc}}$) in Minkowski space (for a gravitational field of nonvanishing field strength it needs the local form of (N$_{+}$), a premise argued at leading order); and in all three a jointly KMS, even state, quasi-free by hypothesis: the joint KMS part of (Q). (Q)'s weight clause, the affine ladder $\kappa(k + \tfrac12)$ of its version 3 ($\kappa k$ in version 2), is not given: the continuum carries every weight. Remark (a comparison, not used). The period $2\pi$ agrees with the period at which the boost's Euclidean section is regular (Proposition \claimref{singularity.euclidean-section-period}, Theorem \claimref{synthesis.local-boost-temperature-is-local-regularity}). The rotation $x = re^{-I_{4}\pi}$ used in (ii) is the imaginary rapidity $\pi$ that carries the half-generator to the other half; the strip $0 < \theta < \pi$ on which the twistor form is definite along the imaginary boost (Theorem \claimref{synthesis.twistor-form-states-only-off-the-real-slice}(iii)) is the $\Delta^{1/2}$ half of the two-point KMS strip $0 < \theta < 2\pi$, with midpoint the $\Delta^{1/4}$ point $\theta = \pi/2$.

\emph{(iv)} For the fields, $\beta_{t}\nu_{a}\beta_{-t}(\psi(u)) = \psi(u + e^{\kappa t}a)$ in the order that conjugates translations by the boost, so $\beta_{t}\nu_{a}\beta_{-t} = \nu_{e^{\kappa t}a}$: the boost dilates $\nu$ at rate $\kappa$. If the boost multiplied each of a generating set of modes by a phase, Theorem \claimref{synthesis.point-spectrum-boost-admits-no-half-sided-translations}(ii) would make $\nu$ the identity, and (iii) of that Theorem would make its unitary implementer trivial on the GNS space; (N$_{+}$) would then have no content (a trivial translation group has the zero generator). So (N$_{+}$) is posed on one-particle spaces where the boost has no eigenvectors spanning the relevant vectors: here $h$ is multiplication by $-\kappa w$ on $L^{2}(\R, dw)$, purely continuous. The stored premises are worded on point-spectrum ladders with phase-acting flows; on the continuum the boost has no eigenvectors, a normalisable mode $f$ has occupation $\ip{f}{A_{+}f} = \int|\hat f(w)|^{2}m(w)\,dw$ in the Mellin transform $\hat f$, and a single mode's algebra is not boost-invariant, so the ladder statements are discretisations of (ii), recovering $m(w_{j})$ only for modes of vanishing Mellin width; no window carries (N$_{+}$). The check records the finite form of the obstruction: on a window with real weights $k_{j}$, $(k_{j} - k_{l})P_{jl} = JcP_{jl}$ with $c$ real and nonzero forces $P = 0$ ($J = e_{2}e_{3}$ the unit of the windows, as in Theorem \claimref{synthesis.point-spectrum-boost-admits-no-half-sided-translations}).

\emph{(v)} Write $\ell = (1, \hat n)$, $|\hat n| = 1$, and $p = (p_{0}, \vec p)$ with $p_{0} \ge |\vec p|$. Then $\ip{\ell}{p} = p_{0} - \hat n\cdot\vec p \ge p_{0} - |\vec p| \ge 0$, with equality only if $\vec p = |\vec p|\hat n$ and $p_{0} = |\vec p|$, that is, $p$ on the ray of $\ell$. So in a representation with spectrum in the closed forward cone the generator $\ip{\ell}{P}$ of the null translations along $\ell$ is positive, and the null-translation group restricted to a null plane meets the positivity half of (N$_{+}$) for every invariant vector state. The check verifies this on 64 rational momenta and seven rational unit directions from Pythagorean quadruples, with equality only on the ray, and that a spacelike translation takes both signs. In the conformal model of Lemma \claimref{mass.chart-form-spin-isomorphism} ($\Cl(2,4)$, $\ip{\nO}{\nInf} = 1$, $X(x) = \nO + x - \tfrac12 Q(x)\nInf$), the rotor $T_{\alpha} = 1 + \tfrac{\alpha}{2}\ell\nInf$ maps $X(u\ell)$ to $X((u + \alpha)\ell)$; the boost $B$ above and the dilation $D = \cosh\tfrac{\lambda}{2} + \sinh\tfrac{\lambda}{2}\,e_{-}e_{+}$ (both bivectors square to $1$) send $X(u\ell)$ to $X(e^{\lambda}u\ell)$ and to $e^{-\lambda}X(e^{\lambda}u\ell)$, the same point, while $B$ fixes $X(e_{2})$ and $D$ moves it; and $BT_{\alpha}\tilde B = DT_{\alpha}\tilde D = T_{e^{\lambda}\alpha}$. This is the relation $\Delta^{I_{4}s}U(a)\Delta^{-I_{4}s} = U(e^{-2\pi s}a)$ of the Borchers form of Fact \claimref{angle.imported-theorems-deliver-their-hypotheses} (the unit of its Hilbert space read as $I_{4}$, numbers in $\mathrm{span}\{1, I_{4}\}$, as for $\mathfrak{h}$), read with the modular group as the boost at $-2\pi s$: the computation of (ii) is an explicit instance of that theorem on the generator, which is not otherwise used.
